\begin{enumerate}
\item $(u+v \mid u-v)=\|u\|^2-\|v\|^2=0$ pour $u$ et $v$ unitaires.
\item\label{ENDO-019-cor-2} Soient $u$ et $v$ des vecteurs unitaires de $E$.
$u+v$ et $u-v$ sont orthogonaux donc $f(u+v)$ et $f(u-v)$ le sont aussi.
Or par linéarité
$$
f(u+v)=f(u)+f(v) \text { et } f(u-v)=f(u)-f(v)$$
de sorte que l'orthogonalité de ces deux vecteurs entraîne
$$
\|f(u)\|=\|f(v)\|
$$
Ainsi les vecteurs unitaires de $E$ sont envoyés par $f$ sur des vecteurs ayant tous la même norme $\alpha \in \mathbb{R}^{+}$.
Montrons qu'alors
$$
\forall\, x \in E,\ \|f(x)\|=\alpha\|x\|
$$
Soit $x \in E$.\\
Si $x=0$ alors on a $f(x)=0$ puis $\|f(x)\|=\alpha\|x\|$.\\
Si $x \neq 0$ alors en introduisant le vecteur unitaire $u=x /\|x\|$, on a $\|f(u)\|=\alpha$ puis $\|f(x)\|=\alpha\|x\|$
\item Si $\alpha=0$ alors $f=\tilde{0}$ et n'importe quel $g \in \mathrm{O}(E)$ convient.\\
Si $\alpha \neq 0$ alors introduisons l'endomorphisme
$$
g=\frac{1}{\alpha} f
$$
La relation obtenue en \textbf{\ref{ENDO-019-cor-2}} assure que $g$ conserve la norme et donc $g \in \mathrm{O}(E)$ ce qui permet de conclure.                                                                                                                    \end{enumerate}
